\section{Création de la base de données}

Après avoir nettoyé les textes, il a fallu créer une base de données pour stocker les informations relatives aux romans. J'ai choisi d'utiliser le système de gestion de base de données relationnelle \textsl{MySQL}, car il permet d'organiser les données sous forme de tables et de conserver des relations claires entre les auteurs, les romans et les unités textuelles extraites.

Je vais donc présenter la structure de cette base de données ainsi que les différentes tables qui la composent. Chaque table a un rôle précis : certaines servent à stocker les informations bibliographiques, tandis que d'autres permettent de conserver les textes, les segments ou encore les phrases extraites pour l'analyse.

J'ai choisi de présenter ce code directement dans le corps du mémoire\footnote{Je précise également que tous les codes utilisés pour la réalisation de ce mémoire sont disponibles sur mon GitHub.} et non uniquement en annexe, afin de rendre plus claire la manière dont les données ont été organisées avant l'analyse. 

\begin{listing}[htbp]
\begin{minted}[
    fontsize=\scriptsize,
    breaklines=true,
    frame=lines,
    linenos=true,
    numbersep=5pt,
    autogobble
]{sql}
CREATE DATABASE romans_nlp;
USE romans_nlp;

-- Creation des tables pour la base de donnees des romans et des auteurs
CREATE TABLE auteurs (
    id_auteur INT PRIMARY KEY AUTO_INCREMENT,
    nom VARCHAR(100) NOT NULL,
    prenom VARCHAR(100)
);

-- Table pour stocker les romans, avec leur titre et leur annee de publication
CREATE TABLE romans (
    id_roman INT PRIMARY KEY AUTO_INCREMENT,
    titre VARCHAR(255) NOT NULL,
    annee_publication INT,
    texte_brut LONGTEXT
);

-- Table d'association pour gerer la relation plusieurs-a-plusieurs entre auteurs et romans
CREATE TABLE auteurs_romans (
    id_auteur INT NOT NULL,
    id_roman INT NOT NULL,
    PRIMARY KEY (id_auteur, id_roman),
    FOREIGN KEY (id_auteur) REFERENCES auteurs(id_auteur),
    FOREIGN KEY (id_roman) REFERENCES romans(id_roman)
);

-- Table pour stocker les segments extraits des romans, avec leur ordre de segmentation
CREATE TABLE segments (
    id_segment INT PRIMARY KEY AUTO_INCREMENT,
    id_roman INT NOT NULL,
    ordre_segment INT NOT NULL,
    texte_segment LONGTEXT NOT NULL,
    FOREIGN KEY (id_roman) REFERENCES romans(id_roman)
);

/*
Apres reflexion, il est egalement possible d'ajouter les dates
de naissance et de mort des auteurs afin d'enrichir la base de donnees.
*/
ALTER TABLE auteurs
ADD date_naissance INT;

ALTER TABLE auteurs
ADD date_mort INT;

-- Table pour stocker les phrases extraites des romans, avec leur motif d'extraction
CREATE TABLE phrase_roman (
    id_phrase INT PRIMARY KEY AUTO_INCREMENT,
    id_roman INT NOT NULL,
    ordre_phrase INT NOT NULL,
    texte_phrase TEXT NOT NULL,
    motif_phrase TEXT NOT NULL,
    FOREIGN KEY (id_roman) REFERENCES romans(id_roman)
);
\end{minted}
\caption{Création de la base de données romans\_nlp}
\end{listing}

Ce code permet de créer la base de données utilisée pour stocker les romans du corpus ainsi que les informations nécessaires à leur analyse. La base, nommée \texttt{romans\_nlp}, est d'abord créée puis sélectionnée afin que toutes les tables suivantes y soient enregistrées.

La première table, \texttt{auteurs}, sert à stocker les informations relatives aux auteurs. Chaque auteur reçoit un identifiant unique, auquel sont associés son nom et, lorsque l'information est disponible, son prénom. Cette table est ensuite complétée par deux colonnes supplémentaires, \texttt{date\_naissance} et \texttt{date\_mort}, afin d'enrichir la description des auteurs et de permettre d'éventuelles analyses chronologiques.

La table \texttt{romans} contient les informations principales sur les œuvres du corpus. Elle associe à chaque roman un identifiant unique, un titre, une année de publication ainsi que le texte brut du roman. Le choix du type \texttt{LONGTEXT} pour le champ \texttt{texte\_brut} permet de stocker des textes longs, ce qui est nécessaire dans le cas de romans complets.

La table \texttt{auteurs\_romans} permet de faire le lien entre les auteurs et les romans. Elle sert à gérer une relation de type plusieurs-à-plusieurs : un auteur peut être associé à plusieurs romans, et un roman peut éventuellement être associé à plusieurs auteurs. Les clés étrangères \texttt{id\_auteur} et \texttt{id\_roman} permettent de relier cette table aux tables \texttt{auteurs} et \texttt{romans}.

La table \texttt{segments} sert à stocker les segments extraits des romans. Chaque segment est rattaché à un roman grâce à la clé étrangère \texttt{id\_roman}. Le champ \texttt{ordre\_segment} permet de conserver l'ordre d'apparition des segments dans le texte d'origine. Cela est important, car même si le texte est découpé pour les traitements, il reste possible de retrouver sa progression initiale.

Enfin, la table \texttt{phrase\_roman} permet de stocker les phrases extraites des romans. Chaque phrase est associée au roman dont elle provient, à son ordre d'apparition, à son contenu textuel et au motif qui a permis son extraction. Ce dernier champ est important, car il permet de garder une trace de la méthode utilisée pour sélectionner les phrases et de relier chaque extraction à un critère précis.

L'ensemble de cette structure permet donc de conserver à la fois les textes complets, leurs découpages en segments, les phrases extraites et les informations bibliographiques liées aux auteurs et aux œuvres. Elle constitue ainsi une base organisée pour les traitements de fouille de texte et d'analyse automatique du corpus.
